\documentclass[11pt,a4paper]{article}
\usepackage[T1]{fontenc}
\usepackage[utf8]{inputenc}
\usepackage{lmodern}
\usepackage{amsmath,amssymb,amsthm}
\usepackage{booktabs}
\usepackage[margin=1.1in]{geometry}
\usepackage{microtype}
\usepackage[hidelinks]{hyperref}

\theoremstyle{plain}
\newtheorem{theorem}{Theorem}[section]
\newtheorem{proposition}[theorem]{Proposition}
\newtheorem{lemma}[theorem]{Lemma}
\newtheorem{corollary}[theorem]{Corollary}
\theoremstyle{definition}
\newtheorem{definition}[theorem]{Definition}
\newtheorem{example}[theorem]{Example}
\theoremstyle{remark}
\newtheorem{remark}[theorem]{Remark}
\newtheorem{observation}[theorem]{Computational observation}

\newcommand{\HH}{\mathbb{H}^2}
\newcommand{\ZZ}{\mathbb{Z}}
\newcommand{\QQ}{\mathbb{Q}}
\newcommand{\RR}{\mathbb{R}}
\newcommand{\vc}[1]{(#1)}
\DeclareMathOperator{\Res}{Res}
\DeclareMathOperator{\area}{area}

\title{Exact algebraic structure of hyperbolic hybrid tilings}
\author{Marek \v{C}trn\'{a}ct \and Alessandro Longo}
\date{Draft of \today}

\begin{document}
\maketitle

\begin{abstract}
A \emph{hybrid} system in the hyperbolic plane is a collection of vertex
combinations of regular polygons that close at a common edge length, so that they
may occur in a single tiling. Catalogues of such systems have so far been built
numerically: the shared edge lengths are computed by iteration to some fifty
decimal places, and the coincidences are believed but not proved. We show that
these edge lengths are algebraic numbers, that a defining polynomial for each is
computable by elimination, and that the question of whether two vertex
combinations share an edge is decidable exactly. Applying the procedure we give
closed forms for eight catalogued edges, among them
$\cosh s=\varphi$, $\cosh s=(3+8\sqrt2)/7$ and $s=\log 3$, and we prove exact
every identity in those systems. The apparatus is insensitive to the sign of the
curvature, so it treats spherical and Euclidean systems at the same time. For the system at $\cosh s=\varphi$ we then give
a complete additive description: all twelve polygons occurring at integer
multiples of the base edge have interior angles lying in a rank-three lattice
spanned by $\pi$, $\arccos(1/\varphi)$ and $\arccos(-1/\varphi^{2})$, so that
enumerating the vertex combinations reduces to three linear Diophantine
equations. Finally we classify the systems in which an equilateral triangle at
$r$ times the base edge has $r$ times the area, obtaining
$\cosh s=1+\sqrt[3]{2}+\sqrt[3]{4}$, $(5+\sqrt{33})/4$, $(3+\sqrt{33})/4$ and
$\varphi$ for $r=2,3,4,7$, and showing that these are the only values of $r$ for
which further polygons occur.
\end{abstract}

\section{Introduction}

Let $\HH$ denote the hyperbolic plane of curvature $-1$. A regular $n$-gon with
side length $s$ exists in $\HH$ for every $n\ge 3$ and every $s>0$, and its
interior angle $\theta_n(s)$ decreases from the Euclidean value
$(1-2/n)\pi$ towards $0$ as $s$ grows. The apeirogon, which we write $n=\infty$,
is the limiting case: a regular polygon with infinitely many sides of length $s$
inscribed in a horocycle.

Because the angle depends on the side, polygons of different shapes can be made to
fit around a vertex by choosing the side length appropriately. This is the
mechanism behind what the first author's working notes call \emph{hybrid}
systems: a set of vertex combinations that all close at one and the same edge
length. Such a set is the raw material of a tiling in which several distinct
regular polygons occur, all with a common side.

The catalogues of these systems have been assembled experimentally. For each
candidate vertex combination one solves a transcendental equation numerically for
the edge, records the result to high precision, and groups combinations whose
recorded edges agree. Agreement to fifty decimal places is persuasive, and a
mismatch at that precision would be strange, but it is not a proof. Recording
this as an open problem, those notes ask for the identities to be established
exactly.

That is the first thing we do here. The key observation is elementary: the
closure condition at a vertex, which as stated is a transcendental equation in
the edge, becomes a \emph{multiplicative} relation among algebraic numbers on the
unit circle once one passes to half-angles (Theorem~\ref{thm:closure}). From this
the edge is seen to be algebraic, an explicit defining polynomial is computable by
resultants, and the question of whether two combinations share an edge becomes a
greatest common divisor computation (Theorem~\ref{thm:decide}). Section~\ref{sec:table}
carries this out for eight systems of the catalogue.

The remainder of the paper studies one system in detail, the one at
$\cosh s=\varphi$ where $\varphi=(1+\sqrt5)/2$, together with the polygons that
occur at integer multiples of its base edge. Section~\ref{sec:lattice} shows that
all twelve such polygons have interior angles in a lattice of rank three, which
reduces the enumeration of their vertex combinations to linear algebra over
$\ZZ$. Section~\ref{sec:star} isolates the geometric construction behind two of
the coincidences and determines exactly when it applies.
Section~\ref{sec:rfamily} classifies a one-parameter family of systems singled
out by an area condition.

Throughout we distinguish carefully between statements we prove and statements
established by computation over a bounded search range; the latter are labelled
as computational observations.

\section{Preliminaries}

\begin{lemma}\label{lem:basic}
Let $P$ be a regular $n$-gon in $\HH$ with side length $s$ and interior angle
$\theta$, where $3\le n\le\infty$. Write $c=\cosh(s/2)$ and $c_n=\cos(\pi/n)$,
with $c_\infty=1$. Then
\begin{equation}\label{eq:basic}
  \sin\frac{\theta}{2}=\frac{c_n}{c},
  \qquad
  \cos\frac{\theta}{2}=\frac{\sqrt{c^{2}-c_n^{2}}}{c}.
\end{equation}
The polygon exists precisely when $c>c_n$, and then $\theta\in(0,\pi)$.
\end{lemma}

\begin{proof}
Join the centre of $P$ to a vertex and to the midpoint of an adjacent side. This
produces a right triangle with legs meeting at the centre in the angle $\pi/n$,
at the vertex in the angle $\theta/2$, and with the side opposite the central
angle equal to $s/2$. The hyperbolic right-triangle relation
$\cosh a=\cos A/\sin B$, where $a$ is the side opposite $A$ and $B$ is the third
angle, gives $\cosh(s/2)=\cos(\pi/n)/\sin(\theta/2)$, which is the first identity.
Since $\theta\in(0,\pi)$ we have $\theta/2\in(0,\pi/2)$, so the cosine is
positive and the second identity follows. The apeirogon is the limit $n\to\infty$,
for which $c_n\to1$.
\end{proof}

\begin{remark}[Curvature is a sign, not a case]\label{rem:curvature}
Nothing below uses $c>1$. In all three geometries of constant curvature the same
relation holds,
\[
  c=\frac{\cos(\pi/n)}{\sin(\theta/2)},
  \qquad
  c=\begin{cases}
      \cos(s/2) & \text{on the sphere},\\
      1         & \text{in the Euclidean plane},\\
      \cosh(s/2)& \text{in }\HH,
    \end{cases}
\]
the spherical case being the substitution $s\mapsto is$. Thus $c$ alone carries
the geometry: a solution of the polynomial $P_m$ of Theorem~\ref{thm:algebraic}
with $c<1$ is a spherical vertex, $c=1$ a Euclidean one, and $c>1$ hyperbolic.
For instance the cube $\{4,3\}$ has $c=\cos(\pi/4)/\sin(\pi/3)=\sqrt{2/3}$, and
its edge is $2\arccos c=70.5288\dots^{\circ}$. We phrase everything for $\HH$
because that is where the systems are plentiful, but every statement that does not
mention $\cosh$ applies verbatim to the other two geometries.
\end{remark}

\begin{remark}[Shortchord, and cyclic polygons]\label{rem:shortchord}
The constant $c_n=\cos(\pi/n)$ is half the \emph{shortchord} of a Euclidean
regular $n$-gon of side $1$, that is half its shortest diagonal. With that
reading, \eqref{eq:closure} says
$\sum_i 2\arcsin\bigl(a_i/(2R)\bigr)=2\pi$ with $a_i=2c_{n_i}$ and $R=c$, which
is exactly the condition that a Euclidean polygon with sides $a_i$ be inscribed
in a circle of radius $R$ containing its centre. So $P_m$ is also a defining
polynomial for the circumradius of a cyclic polygon with prescribed sides. That
problem is not open: Robbins~\cite{robbins} obtained the analogous area
polynomial in 1994 and Pech~\cite{pech} the circumradius polynomial in 2006,
both of degree $7$ in the pentagon and heptagon cases, with later work by
Varfolomeev and others; explicit general formulas beyond small $n$ remain hard.
The elimination of Theorem~\ref{thm:algebraic} is an independent route to those
polynomials, and conversely their degree bounds are the natural place to look for
the height bound requested in Section~\ref{sec:limits}.
\end{remark}

For fixed $n$ the map $c\mapsto\theta_n(c)=2\arcsin(c_n/c)$ is strictly
decreasing on $(c_n,\infty)$, with $\theta_n(c)\to\pi-2\pi/n$ as $c\downarrow c_n$
and $\theta_n(c)\to0$ as $c\to\infty$.

\begin{definition}[Vertex combination and edge function]
A \emph{vertex combination} is a function $m$ from $\{3,4,5,\dots\}\cup\{\infty\}$
to $\ZZ_{\ge0}$ with finite support, written multiplicatively as
$\vc{n_1^{m_1}.n_2^{m_2}\cdots}$. Its \emph{degree} is $M=\sum_n m_n$. We say $m$
\emph{closes} at $c$ if
\begin{equation}\label{eq:closure}
  \sum_n m_n\,\theta_n(c)=2\pi .
\end{equation}
\end{definition}

\begin{lemma}\label{lem:unique}
A vertex combination closes at no more than one value of $c$.
\end{lemma}

\begin{proof}
On the common domain $c>\max\{c_n: m_n>0\}$ each $\theta_n$ is strictly
decreasing, hence so is the finite positive combination
$\sum_n m_n\theta_n$. A strictly monotone function attains the value $2\pi$ at
most once.
\end{proof}

When the value exists we write $e(m)=c$ for it, and call it the \emph{edge
function} of $m$; we shall also refer to $s=2\operatorname{arcosh}c$ as the edge.
Two vertex combinations $m,m'$ with $e(m)=e(m')$ are said to satisfy a
\emph{hybrid identity}; a maximal set of combinations pairwise satisfying one is a
\emph{hybrid system}. Only combinations in a common system can occur in a single
tiling by regular polygons of one side length.

\section{The closure condition is a multiplicative relation}\label{sec:decide}

Fix a vertex combination $m$ with support $S$ and degree $M$. For $n\in S$
introduce the auxiliary quantity $u_n=\sqrt{c^{2}-c_n^{2}}$, so that by
Lemma~\ref{lem:basic}
\[
  z_n \;=\; e^{i\theta_n/2}\;=\;\frac{u_n+i\,c_n}{c}.
\]

\begin{theorem}\label{thm:closure}
With the notation above, $m$ closes at $c$ if and only if
\begin{equation}\label{eq:mult}
  \prod_{n\in S}\bigl(u_n+i\,c_n\bigr)^{m_n}\;=\;-\,c^{M},
\end{equation}
and \eqref{eq:mult} implies, more weakly,
$\sum_n m_n\theta_n\equiv2\pi \pmod{4\pi}$.
\end{theorem}

\begin{proof}
Condition \eqref{eq:closure} is equivalent to $\sum_n m_n\theta_n/2=\pi$, hence to
$\prod_n z_n^{m_n}=e^{i\pi}=-1$. Multiplying by $c^{M}$ gives \eqref{eq:mult}.
Conversely \eqref{eq:mult} gives $\prod_n z_n^{m_n}=-1$, that is
$\sum_n m_n\theta_n/2\equiv\pi\pmod{2\pi}$, which is the stated congruence.
\end{proof}

The point of \eqref{eq:mult} is that its two sides are polynomials in $c$, the
$u_n$, and the algebraic numbers $c_n$, whereas \eqref{eq:closure} is
transcendental. Separating real and imaginary parts, the imaginary part of the
left side of \eqref{eq:mult} must vanish; that is a single polynomial condition.

\begin{theorem}\label{thm:algebraic}
Let $m$ be a vertex combination that closes. Then $e(m)$ is an algebraic number,
and a nonzero polynomial $P_m\in\QQ(\{c_n\}_{n\in S})[X]$ with $P_m(e(m))=0$ is
computable by the following procedure.
\begin{enumerate}
\item Expand $W=\prod_{n\in S}(U_n+i\,c_n)^{m_n}$ as a polynomial in independent
      indeterminates $U_n$.
\item Reduce every $U_n^{2}$ to $X^{2}-c_n^{2}$, so that each $U_n$ occurs to
      degree at most one.
\item Let $E=\operatorname{Im}W$, a polynomial in $X$ and the $U_n$.
\item Successively eliminate the $U_n$ by
      $E\mapsto\Res_{U_n}\!\bigl(E,\;U_n^{2}-(X^{2}-c_n^{2})\bigr)$.
\end{enumerate}
The result $P_m$ is a polynomial in $X$ alone, and it is even in $X$.
\end{theorem}

\begin{proof}
By Theorem~\ref{thm:closure} the number $c=e(m)$, together with the positive
square roots $u_n=\sqrt{c^{2}-c_n^{2}}$, satisfies $\operatorname{Im}W=0$ after
substituting $X=c$, $U_n=u_n$. Steps~2 and~4 are the standard elimination of the
$u_n$: the resultant of $E$ with the minimal polynomial of $u_n$ over
$\QQ(\{c_j\},X)$ lies in the ideal generated by them, so it vanishes at
$X=c$, and it no longer involves $U_n$. After all eliminations $P_m(c)=0$.
Nonvanishing of $P_m$ as a polynomial is checked in the course of the
computation. Since $\theta_n$ depends on $c$ only through $c^{2}$ in
\eqref{eq:basic} after the substitution $u_n^{2}=c^{2}-c_n^{2}$, the surviving
polynomial is even in $X$; we write $P_m(X)=F_m(X^{2})$.
\end{proof}

\begin{remark}
The root set of $P_m$ is in general larger than $\{e(m)\}$. Two sources of
spurious roots are visible in the proof: the congruence in
Theorem~\ref{thm:closure} admits $\sum m_n\theta_n\in\{2\pi,6\pi,\dots\}$, and
the elimination step forgets the sign conditions $u_n>0$. This does not affect
what follows, because we always identify the relevant root by numerical
isolation.
\end{remark}

\begin{theorem}[Decision procedure]\label{thm:decide}
Let $m,m'$ be vertex combinations that close, with $P_m,P_{m'}$ as above, and let
$G=\gcd(P_m,P_{m'})$, computed exactly. Suppose $I\subset\RR$ is an interval such
that
\begin{enumerate}
\item $P_m$ has exactly one root in $I$, and $P_{m'}$ has exactly one root in
      $I$;
\item $e(m)\in I$ and $e(m')\in I$;
\item $G$ has a root in $I$.
\end{enumerate}
Then $e(m)=e(m')$. If instead $G$ has no root in $I$, then $e(m)\ne e(m')$.
\end{theorem}

\begin{proof}
Let $\alpha\in I$ be a root of $G$. Then $\alpha$ is a root of $P_m$, and $e(m)$
is a root of $P_m$ lying in $I$; by hypothesis~(1) they coincide, so
$\alpha=e(m)$. The same argument applied to $P_{m'}$ gives $\alpha=e(m')$. Hence
$e(m)=e(m')$. If $G$ has no root in $I$ then $e(m)$, being in $I$, is not a
common root, so $e(m)\ne e(m')$.
\end{proof}

All three hypotheses are effectively checkable: root counting in an interval is
Sturm's theorem over $\QQ$, or root isolation over the coefficient field in
general, and $e(m)$ is located to any desired precision by
Lemma~\ref{lem:unique} together with bisection. The procedure therefore decides
hybrid identities. This settles, in principle, the question of whether the
numerically observed coincidences are exact; the remaining difficulty is
computational, and we return to it in Section~\ref{sec:limits}.

\section{Exact edges for eight catalogued systems}\label{sec:table}

We have run the procedure of Theorems~\ref{thm:algebraic} and~\ref{thm:decide} on
part of the catalogue. Table~\ref{tab:edges} records the outcome. In each case
the polynomials $F_m$ were computed symbolically, the greatest common divisor was
taken exactly, and the resulting value of $\cosh s=2c^{2}-1$ was solved in closed
form, where the degree permits, and then verified against every member
combination of the system to eighty decimal digits. The quartic at $1.1555$ is
irreducible over $\QQ$; we give its minimal polynomial instead of a radical
expression.

\begin{table}[t]
\centering
\begin{tabular}{@{}lll@{}}
\toprule
$s$ (as catalogued) & vertex combinations & $\cosh s$ exactly \\
\midrule
$0.6196\dots$ & $\vc{3^{2}.5.\infty}$, $\vc{3.5^{3}}$
  & $\dfrac{2+5\sqrt5}{11}$ \\[2mm]
$1.0612\dots$ & $\vc{3^{4}.4^{2}}$, $\vc{3^{2}.4.5^{2}}$, $\vc{3.4.10.20}$, $\vc{5^{4}}$
  & $\varphi=\dfrac{1+\sqrt5}{2}$ \\[2mm]
$1.0986\dots$ & $\vc{4^{2}.5.\infty}$, $\vc{\infty^{3}}$
  & $\dfrac53$, \; so $s=\log 3$ \\[2mm]
$1.1555\dots$ & $\vc{3.5.12^{2}}$, $\vc{4.5^{2}.12}$
  & root of $191y^{4}-424y^{3}+10y^{2}+328y-121$ \\[2mm]
$1.2537\dots$ & $\vc{4^{5}}$, $\vc{4^{2}.10.30}$
  & $\dfrac{5+2\sqrt5}{5}=1+\dfrac{2}{\sqrt5}$ \\[2mm]
$1.3424\dots$ & $\vc{3^{3}.\infty^{2}}$, $\vc{3.8.24.\infty}$
  & $\dfrac{3+8\sqrt2}{7}$ \\[2mm]
$1.6628\dots$ & $\vc{3^{4}.\infty^{2}}$, $\vc{3^{2}.12^{2}.\infty}$, $\vc{12^{4}}$
  & $1+\sqrt3$ \\[2mm]
$1.7627\dots$ & $\vc{4^{6}}$, $\vc{4^{3}.\infty^{2}}$, $\vc{\infty^{4}}$
  & $3$, \; so $s=\log(3+2\sqrt2)$ \\
\bottomrule
\end{tabular}
\caption{Exact edge lengths. Every hybrid identity in these eight systems is a
theorem, by Theorem~\ref{thm:decide}.}
\label{tab:edges}
\end{table}

\begin{example}
Take $\vc{3^{4}.4^{2}}$, four triangles and two squares. Here
$c_3=\tfrac12$, $c_4=\tfrac{\sqrt2}{2}$, $M=6$, and the elimination of
Theorem~\ref{thm:algebraic} yields $P(X)=4X^{4}-6X^{2}+1$, so
$c^{2}=(3+\sqrt5)/4$ and
\[
  \cosh s=2c^{2}-1=\frac{1+\sqrt5}{2}=\varphi .
\]
The same value is returned by $\vc{5^{4}}$, by $\vc{3^{2}.4.5^{2}}$ and by
$\vc{3.4.10.20}$, and in each case the greatest common divisor of the two
polynomials is of positive degree, so the identity is exact.
\end{example}

\begin{remark}
Two of the eight edges satisfy $\cosh s\in\QQ$, whence $e^{s}$ is an algebraic
unit and $s$ is the logarithm of one: $s=\log3$ and $s=\log(3+2\sqrt2)$. We note
in passing that $s$ itself is never algebraic here. Indeed if $s\ne0$ were
algebraic then $\cosh s$ would be transcendental by the Lindemann--Weierstrass
theorem, contradicting Theorem~\ref{thm:algebraic}. This is consistent with, and
a special case of, the theorem of Calcut~\cite{calcutrational} that a hyperbolic
triangle with all angles rational multiples of $\pi$ has transcendental sides.
\end{remark}

\section{The system at \texorpdfstring{$\cosh s=\varphi$}{cosh s = phi}}\label{sec:lattice}

We now fix the system of Table~\ref{tab:edges} with $\cosh s=\varphi$ and study
the polygons occurring at integer multiples of its base edge. Following the
catalogue's notation we write a letter for the multiplier and a number for the
number of sides, so that $\mathrm{A}3$ is a triangle of side $s$, $\mathrm{B}4$ a
square of side $2s$, $\mathrm{G}3$ a triangle of side $7s$, and
$\mathrm{D}\infty$ an apeirogon of side $4s$.

\begin{lemma}\label{lem:base}
The base edge satisfies $\cosh s=\varphi$, and it is the edge of the regular
tiling $\{5,4\}$, and of $\{5,6\}$, $\{3,10\}$ and $\{10,5\}$ at twice its
length. Moreover
\[
  \cosh\frac{s}{2}=\frac{\varphi}{\sqrt2},\quad
  \cosh s=\varphi,\quad
  \cosh 2s=\varphi^{3},\quad
  \cosh\frac{7s}{2}=\frac{\varphi^{7}}{\sqrt2}.
\]
\end{lemma}

\begin{proof}
For $\{5,4\}$, Lemma~\ref{lem:basic} with $n=5$ and $\theta=\pi/2$ gives
$\cosh(s/2)=\cos(\pi/5)/\sin(\pi/4)=\varphi/\sqrt2$, whence
$\cosh s=2\varphi^{2}/2-1=\varphi^{2}-1=\varphi$. For $\{5,6\}$ one finds
$\cosh(s'/2)=\cos(\pi/5)/\sin(\pi/6)=\varphi=\cosh s$, so $s'=2s$; the
computation for $\{3,10\}$ and $\{10,5\}$ is identical, with
$\cos(\pi/3)/\sin(\pi/10)=\varphi$ and $\cos(\pi/10)/\sin(\pi/5)=\varphi$.
Finally $\cosh(ks)=T_k(\cosh s)=T_k(\varphi)$ for the Chebyshev polynomial
$T_k$, and $2\cosh^{2}(ks/2)=T_k(\varphi)+1$; for $k=7$ one computes
$T_7(\varphi)+1=377\varphi+233=\varphi^{14}$, using $F_{14}=377$ and $F_{13}=233$.
\end{proof}

\begin{remark}
Three regular tilings therefore share the edge $2s$: $\{5,6\}$, $\{3,10\}$ and
$\{10,5\}$, that is the vertex combinations $\vc{5^{6}}$, $\vc{3^{10}}$ and
$\vc{10^{5}}$. This is the only coincidence of three regular tilings at one edge
known to us, and Table~\ref{tab:poly} exhibits it as the three polygons
$\mathrm{B}5$, $\mathrm{B}3$, $\mathrm{B}10$ whose angles are $\pi/3$, $\pi/5$
and $2\pi/5$.
\end{remark}

\begin{remark}
The full symmetry group of $\{5,4\}$ is the triangle group $(2,4,5)$, which by
Takeuchi's classification~\cite{takeuchi} is arithmetic with invariant trace
field $\QQ(\sqrt5)$. That is the structural reason the quantities $\cosh(ks)$
are algebraic integers of $\ZZ[\varphi]$.
\end{remark}

\begin{theorem}\label{thm:lattice}
Put $a=\arccos(1/\varphi)$ and $t=\arccos(-1/\varphi^{2})$. The twelve polygons
listed in Table~\ref{tab:poly} have the interior angles shown there. In
particular every one of them lies in the additive group
$\QQ\pi+\ZZ a+\ZZ t$.
\end{theorem}

\begin{table}[t]
\centering
\begin{tabular}{@{}lccll@{}}
\toprule
symbol & $n$ & edge & interior angle & $\cos\theta$\\
\midrule
$\mathrm{A}3$   & $3$      & $s$  & $a$            & $1/\varphi$\\
$\mathrm{A}4$   & $4$      & $s$  & $\pi-2a$       & $1/\varphi^{3}=\sqrt5-2$\\
$\mathrm{A}5$   & $5$      & $s$  & $\pi/2$        & $0$\\
$\mathrm{A}10$  & $10$     & $s$  & $t$            & $-1/\varphi^{2}$\\
$\mathrm{A}20$  & $20$     & $s$  & $\pi+a-t$      & root of $x^{4}+8x^{3}+4x^{2}-8x-4$\\
$\mathrm{B}3$   & $3$      & $2s$ & $\pi/5$        & $(1+\sqrt5)/4$\\
$\mathrm{B}4$   & $4$      & $2s$ & $a$            & $1/\varphi$\\
$\mathrm{B}5$   & $5$      & $2s$ & $\pi/3$        & $1/2$\\
$\mathrm{B}10$  & $10$     & $2s$ & $2\pi/5$       & $(\sqrt5-1)/4$\\
$\mathrm{B}\infty$ & $\infty$ & $2s$ & $\pi-2a$    & $1/\varphi^{3}$\\
$\mathrm{D}\infty$ & $\infty$ & $4s$ & $4a-\pi$    & $8\sqrt5-17$\\
$\mathrm{G}3$   & $3$      & $7s$ & $7a-2\pi$      & $1-\varphi^{-14}$\\
\bottomrule
\end{tabular}
\caption{The twelve polygons of the system, and their angles in the lattice.}
\label{tab:poly}
\end{table}

\begin{proof}[Proof of the first two rows, the rest being similar]
By Lemma~\ref{lem:basic} and Lemma~\ref{lem:base},
$\sin(\theta_3/2)=\tfrac12\cdot\sqrt2/\varphi=1/(\sqrt2\varphi)$ and, using
$2\varphi+1=\varphi^{3}$,
$\cos(\theta_3/2)=\sqrt{\varphi^{3}}/(\sqrt2\varphi)=\sqrt{\varphi}/\sqrt2$.
Hence $\cos\theta_3=1-2\sin^{2}(\theta_3/2)=1-1/\varphi^{2}=1/\varphi$, which is
the assertion $\theta_3=a$. Likewise $\sin(\theta_4/2)=1/\varphi$ and
$\cos(\theta_4/2)=1/\sqrt{\varphi}$, so with
$\zeta_3=e^{i\theta_3/2}=(\varphi^{3/2}+i)/(\sqrt2\varphi)$ and
$\zeta_4=e^{i\theta_4/2}=(\sqrt{\varphi}+i)/\varphi$ we get
\[
  \zeta_3^{2}\zeta_4
  =\frac{1+i\sqrt{\varphi}}{\varphi}\cdot\frac{\sqrt{\varphi}+i}{\varphi}
  =\frac{i(1+\varphi)}{\varphi^{2}}=i ,
\]
that is $\theta_3+\theta_4/2=\pi/2$, or $2\theta_3+\theta_4=\pi$.
\end{proof}

\begin{corollary}\label{cor:g3}
$\cos\mathrm{G}3=L_7/\varphi^{7}=1-\varphi^{-14}=(377\sqrt5-841)/2$, where
$L_7=29$ is a Lucas number.
\end{corollary}

\begin{proof}
Writing $\omega=e^{ia}=(1+i\sqrt\varphi)/\varphi$, which has modulus one because
$1+\varphi=\varphi^{2}$, one has
$\cos(ka)=\operatorname{Re}(1+i\sqrt\varphi)^{k}/\varphi^{k}$.
Expanding, $\operatorname{Re}(1+i\sqrt\varphi)^{7}=1-21\varphi+35\varphi^{2}-7\varphi^{3}=29$,
using $\varphi^{2}=\varphi+1$ and $\varphi^{3}=2\varphi+1$. Since
$\mathrm{G}3=7a-2\pi$ we get $\cos\mathrm{G}3=\cos7a=29/\varphi^{7}$, and
$29=L_7=\varphi^{7}-\varphi^{-7}$ gives the stated forms.
\end{proof}

\begin{observation}\label{obs:indep}
Writing $\operatorname{Re}(1+i\sqrt\varphi)^{k}=p_k+q_k\varphi$ with
$p_k,q_k\in\ZZ$, the recursion $p_k=2p_{k-1}-(p_{k-2}+q_{k-2})$,
$q_k=2q_{k-1}-(p_{k-2}+2q_{k-2})$ gives $q_k=0$ for $k\le20000$ only at $k=1$ and
$k=7$. These are exactly the two multiples of the base edge at which the system
carries an equilateral triangle.
\end{observation}

\begin{observation}\label{obs:rank}
The integer relation algorithm PSLQ, applied at $200$ digits to
$(\pi,a,t)$, finds no relation with coefficients below $10^{15}$. We therefore
treat the lattice of Theorem~\ref{thm:lattice} as having rank three. The margin
matters for what follows only through coefficients of size at most about $900$,
since a vertex combination is bounded by $2\pi$ divided by the smallest angle.
\end{observation}

\subsection{Enumerating the vertex combinations}

Granting Observation~\ref{obs:rank}, closing a vertex becomes linear. Writing
$m(\cdot)$ for multiplicities and clearing denominators, the three coordinates of
$\sum m(P)\theta_P=2\pi$ in the basis $(t,a,\pi)$ read
\begin{align}
  m(\mathrm{A}10)-m(\mathrm{A}20)&=0,\label{eq:t}\\
  m(\mathrm{A}3)+m(\mathrm{B}4)+m(\mathrm{A}20)-2m(\mathrm{A}4)-2m(\mathrm{B}\infty)
  +4m(\mathrm{D}\infty)+7m(\mathrm{G}3)&=0,\label{eq:a}\\
  \begin{split}
  30m(\mathrm{A}4)+15m(\mathrm{A}5)+30m(\mathrm{A}20)+6m(\mathrm{B}3)+10m(\mathrm{B}5)\hspace{6em}&\\
  {}+12m(\mathrm{B}10)+30m(\mathrm{B}\infty)-30m(\mathrm{D}\infty)-60m(\mathrm{G}3)&=60 .
  \end{split}\label{eq:pi}
\end{align}

\begin{observation}\label{obs:106}
The system \eqref{eq:t}--\eqref{eq:pi} has exactly $106$ solutions in
non-negative integers; these fall into $30$ classes if polygons are identified by
angle rather than by shape. The analogous system with right-hand side $\pi$,
describing vertices lying in the interior of a longer edge, has exactly $14$
solutions in $7$ angle classes. Each of the twelve polygons occurs in at least
one solution.
\end{observation}

The count $106$ agrees with the catalogue. Two structural consequences are worth
recording, both immediate from the equations.

\begin{proposition}
In any solution of \eqref{eq:t}--\eqref{eq:pi}, $m(\mathrm{A}10)=m(\mathrm{A}20)$.
Consequently a $20$-gon and a decagon occur at exactly the same vertices, and in
equal numbers.
\end{proposition}

\begin{proof}
Immediate from \eqref{eq:t}.
\end{proof}

\begin{proposition}\label{prop:g3}
In any solution, $m(\mathrm{G}3)\le1$.
\end{proposition}

\begin{proof}
By \eqref{eq:a},
$7m(\mathrm{G}3)\le 2\bigl(m(\mathrm{A}4)+m(\mathrm{B}\infty)\bigr)$. The
polygons $\mathrm{A}4$ and $\mathrm{B}\infty$ both have interior angle
$\pi-2a$, and $5(\pi-2a)>2\pi$ numerically, so at most four of them meet at a
vertex. Hence $7m(\mathrm{G}3)\le8$ and $m(\mathrm{G}3)\le1$.
\end{proof}

Proposition~\ref{prop:g3} is sharp in an instructive way: the bound is $8/7$, so
a single unit of slack separates this system from one in which two large
triangles could meet. In the systems of Section~\ref{sec:rfamily} with
$r\in\{2,3,4\}$ the corresponding coefficient is $r<8/2$ and repetitions do occur.

\section{Augmentation, and when it applies}\label{sec:star}

Two of the coincidences in Table~\ref{tab:poly}, namely
$\mathrm{B}4=\mathrm{A}3$ and $\mathrm{B}\infty=\mathrm{A}4$, deserve separate
comment because only the first has a geometric explanation.

\begin{theorem}[Augmentation]\label{thm:star}
Let $s>0$ and $3\le m\le\infty$, and suppose
$\theta_m(s)+2\theta_3(s)=\pi$. Then $\theta_m(2s)=\theta_3(s)$.
\end{theorem}

\begin{proof}
Write $c=\cosh(s/2)$ and $u=\sin(\theta_3/2)=1/(2c)$ by Lemma~\ref{lem:basic}.
The hypothesis gives $\sin(\theta_m/2)=\cos\theta_3=1-2u^{2}$, hence
$\cos(\pi/m)=c\,(1-2u^{2})=(1-2u^{2})/(2u)$. At the doubled edge
$\cosh(s)=2c^{2}-1=1/(2u^{2})-1$, so
\[
  \sin\frac{\theta_m(2s)}{2}
  =\frac{\cos(\pi/m)}{\cosh s}
  =\frac{1-2u^{2}}{2u}\cdot\frac{2u^{2}}{1-2u^{2}}
  =u=\sin\frac{\theta_3(s)}{2}. \qedhere
\]
\end{proof}

Geometrically: attach a triangle of side $s$ to each side of the $m$-gon. At each
original vertex the angle becomes $\theta_m+2\theta_3=\pi$, so those vertices
straighten away, and what remains is an $m$-gon whose sides consist of two base
edges each and whose angles are $\theta_3$. The hypothesis is exactly the
statement that a vertex of type $\vc{3^{4}.m^{2}}$ closes.

\begin{theorem}[Tripling]\label{thm:triple}
Let $s>0$ and $2\le n\le\infty$, and suppose $\theta_3(s)+\theta_{2n}(s)=\pi$.
Then $\theta_n(3s)=\theta_3(s)$.
\end{theorem}

\begin{proof}
Write $c$ for the parameter of Remark~\ref{rem:curvature} and
$u=\sin(\theta_3/2)=1/(2c)$. The hypothesis gives
$\sin(\theta_{2n}/2)=\cos(\theta_3/2)$, hence
$\cos(\pi/2n)=c\sqrt{1-u^{2}}$ and
\[
  \cos(\pi/n)=2\cos^{2}(\pi/2n)-1=2c^{2}\Bigl(1-\tfrac{1}{4c^{2}}\Bigr)-1
  =\tfrac{4c^{2}-3}{2}.
\]
Since the parameter at the tripled edge is $T_3(c)=4c^{3}-3c=c(4c^{2}-3)$,
\[
  \sin\frac{\theta_n(3s)}{2}=\frac{\cos(\pi/n)}{T_3(c)}
  =\frac{4c^{2}-3}{2\,c\,(4c^{2}-3)}=\frac{1}{2c}=u . \qedhere
\]
\end{proof}

Geometrically the hypothesis says that a vertex of type $\vc{3.2n.3.2n}$ closes.
Attach a triangle of side $s$ to every second side of the $2n$-gon. Each of its
$2n$ vertices then carries $\theta_{2n}+\theta_3=\pi$ and straightens away,
leaving an $n$-gon whose sides are three base edges long and whose angles are
$\theta_3$. By Remark~\ref{rem:curvature} the theorem is not confined to $\HH$:
at $n=2$ it applies on the sphere to the cuboctahedron $\vc{3.4.3.4}$, where
$c=\sqrt3/2$, the edge is $60^{\circ}$, and the conclusion degenerates correctly
to $\cos(\pi/2)=0$.

\begin{theorem}[Doubling criterion]\label{thm:double}
For $3\le n,m\le\infty$ and $s>0$,
\[
  \theta_m(2s)=\theta_n(s)
  \iff
  \frac{\cos(\pi/m)}{\cos(\pi/n)}=\frac{\cosh s}{\cosh(s/2)} .
\]
\end{theorem}

\begin{proof}
By Lemma~\ref{lem:basic}, $\theta_m(2s)=\theta_n(s)$ if and only if
$\cos(\pi/m)/\cosh(s)=\cos(\pi/n)/\cosh(s/2)$.
\end{proof}

\begin{proposition}\label{prop:twoonly}
In the system of Section~\ref{sec:lattice} one has
$\cosh s/\cosh(s/2)=\sqrt2$, and the equation
$\cos(\pi/m)=\sqrt2\cos(\pi/n)$ has exactly two solutions with
$3\le n,m\le\infty$, namely $(n,m)=(3,4)$ and $(n,m)=(4,\infty)$. These give
$\mathrm{B}4=\mathrm{A}3$ and $\mathrm{B}\infty=\mathrm{A}4$ respectively.
\end{proposition}

\begin{proof}
$\cosh s/\cosh(s/2)=\varphi/(\varphi/\sqrt2)=\sqrt2$. Since
$\cos(\pi/m)\le1$ we need $\cos(\pi/n)\le1/\sqrt2$, hence $\pi/n\ge\pi/4$ and
$n\le4$. For $n=3$, $\cos(\pi/m)=\sqrt2/2$ forces $m=4$; for $n=4$,
$\cos(\pi/m)=1$ forces $m=\infty$.
\end{proof}

\begin{proposition}
The hypothesis of Theorem~\ref{thm:star}, with a general $n$-gon in place of the
triangle, is equivalent to the criterion of Theorem~\ref{thm:double} only when
$n=3$.
\end{proposition}

\begin{proof}
Writing $x=\cos(\pi/n)$ and $c=\cosh(s/2)$, the augmentation hypothesis
$\theta_m(s)+2\theta_n(s)=\pi$ gives $\cos(\pi/m)=c-2x^{2}/c$, while the
criterion gives $\cos(\pi/m)=x(2c-1/c)$. Their difference, multiplied by $c$,
factors as $(1-2x)(c^{2}+x)$. The second factor is positive, so the two agree if
and only if $x=1/2$, that is $n=3$.
\end{proof}

Thus $\mathrm{B}4=\mathrm{A}3$ arises from the augmentation picture, while
$\mathrm{B}\infty=\mathrm{A}4$ satisfies the criterion without any such
construction behind it.

\section{Triangles of proportional area}\label{sec:rfamily}

The system of Section~\ref{sec:lattice} contains a triangle at seven times the
base edge whose area is also exactly seven times that of the base triangle. We
classify this phenomenon.

\begin{proposition}\label{prop:rcond}
Let $s>0$, let $r\ge2$ be an integer, and write $\theta(s)$ for the interior
angle of the equilateral triangle of side $s$. The following are equivalent:
\begin{enumerate}
\item $\area\bigl(\triangle(rs)\bigr)=r\cdot\area\bigl(\triangle(s)\bigr)$;
\item $\theta(rs)=r\,\theta(s)-(r-1)\pi/3$.
\end{enumerate}
Moreover for each $r$ there is exactly one such $s$.
\end{proposition}

\begin{proof}
The area of a hyperbolic triangle with angles $\theta$ is $\pi-3\theta$, so (1)
reads $\pi-3\theta(rs)=r(\pi-3\theta(s))$, which rearranges to (2). For
uniqueness, put $f(s)=\area(\triangle(rs))/\area(\triangle(s))$. As $s\to0$ the
geometry is asymptotically Euclidean and areas scale as the square of the side,
so $f(s)\to r^{2}$; as $s\to\infty$ both areas tend to $\pi$ and $f(s)\to1$. The
function $f$ is continuous, and computation shows it to be strictly decreasing,
so it attains the value $r\in(1,r^{2})$ exactly once.
\end{proof}

Condition (2) is again a relation between half-angles at two edge multiples, and
the elimination of Theorem~\ref{thm:algebraic} applies verbatim. Writing
$\zeta=e^{i\theta(s)/2}=(u+\tfrac{i}{2})/c$ and
$\zeta_r=e^{i\theta(rs)/2}=(u_r+\tfrac{i}{2})/T_r(c)$ with $c=\cosh(s/2)$,
$u=\sqrt{c^{2}-\tfrac14}$ and $u_r=\sqrt{T_r(c)^{2}-\tfrac14}$, condition (2)
says
\begin{equation}\label{eq:rcond}
  \zeta^{\,r}=e^{i(r-1)\pi/6}\,\zeta_r ,
\end{equation}
whose imaginary part, after eliminating $u$ and $u_r$ by resultants, is a
polynomial in $c$ that is even, hence a polynomial in $\cosh s=2c^{2}-1$.

\begin{theorem}\label{thm:rvalues}
The edge of Proposition~\ref{prop:rcond} satisfies, for $r=2,3,4,7$
respectively,
\[
  \cosh s\;=\;1+\sqrt[3]{2}+\sqrt[3]{4}\;=\;\frac{1}{\sqrt[3]{2}-1},
  \qquad \frac{5+\sqrt{33}}{4},
  \qquad \frac{3+\sqrt{33}}{4},
  \qquad \varphi,
\]
with respective minimal polynomials over $\QQ$
\[
  y^{3}-3y^{2}-3y-1,\qquad 2y^{2}-5y-1,\qquad 2y^{2}-3y-3,\qquad y^{2}-y-1 .
\]
In particular $\cosh s_{3}-\cosh s_{4}=\tfrac12$, and indeed
$y\mapsto y+\tfrac12$ carries $2y^{2}-5y-1$ to $2y^{2}-3y-3$.
\end{theorem}

\begin{proof}
Carry out the elimination described above for \eqref{eq:rcond}, factor the
resulting polynomial in $\cosh s$ over $\QQ$, and select the irreducible factor
vanishing at the numerical value of $\cosh s$; the four factors are those listed.
Solving each gives the stated closed forms. For $r=2$ the cubic
$y^{3}-3y^{2}-3y-1$ is irreducible over $\QQ$; substituting $y=1+z+z^{2}$ and
reducing modulo $z^{3}=2$ gives $0$, so $1+\sqrt[3]{2}+\sqrt[3]{4}$ is its real
root, and $1+z+z^{2}=(z^{3}-1)/(z-1)=1/(\sqrt[3]{2}-1)$.
\end{proof}

The case $r=7$ is the system of Section~\ref{sec:lattice}: the triangle
$\mathrm{G}3$ has $\theta=7a-2\pi$ and therefore area $7(\pi-3a)$.

\begin{table}[t]
\centering
\begin{tabular}{@{}llll@{}}
\toprule
 & $r=2$ & $r=3$ & $r=4$ \\
\midrule
$\cosh s$ & $1+\sqrt[3]{2}+\sqrt[3]{4}$ & $(5+\sqrt{33})/4$ & $(3+\sqrt{33})/4$\\
rank      & $2$ & $3$ & $4$\\
basis     & $\pi,a$ & $\pi,a,g$ & $\pi,a,b,d$\\
\midrule
          & $\mathrm{A}3=a$          & $\mathrm{A}3=a$            & $\mathrm{A}3=a$\\
          & $\mathrm{A}6=\pi-3a$     & $\mathrm{A}6=g$            & $\mathrm{A}6=\pi-2a$\\
          & $\mathrm{B}3=2a-\pi/3$   & $\mathrm{B}3=\pi-2g$       & $\mathrm{B}6=a$\\
          & $\mathrm{B}\infty=\pi-4a$& $\mathrm{C}3=3a-2\pi/3$    & $\mathrm{B}\infty=b$\\
          &                          & $\mathrm{C}6=\pi/3-a$      & $\mathrm{D}3=4a-\pi$\\
          &                          &                            & $\mathrm{D}\infty=d$\\
          &                          &                            & $\mathrm{F}\infty=\pi-3b-d$\\
\midrule
vertices  & $16$ & $12$ & $14$\\
flat      & $6$  & $5$  & $5$\\
\bottomrule
\end{tabular}
\caption{The three smaller systems of Proposition~\ref{prop:rcond}, reduced. The
letter records the edge multiple, so $\mathrm{C}6$ is a hexagon of side $3s$.
Here $a=\theta_3(s)$ throughout, while $g$, $b$, $d$ are the extra generators
carried respectively by $\mathrm{A}6$, $\mathrm{B}\infty$ and $\mathrm{D}\infty$.}
\label{tab:rsys}
\end{table}

The polygons themselves are listed in Table~\ref{tab:rsys}; the system $r=7$ is
Table~\ref{tab:poly}. In each case the base tiling recorded in the catalogue
appears among the vertex combinations: $\vc{\mathrm{A}3^{6}.\mathrm{A}6^{2}}$ at
$r=2$, $\vc{\mathrm{C}3^{6}.\mathrm{C}6^{18}}$ at $r=3$ and
$\vc{\mathrm{A}3^{4}.\mathrm{A}6^{2}}$ at $r=4$.

\begin{observation}\label{obs:rich}
Reducing each of these systems as in Section~\ref{sec:lattice}, and scanning all
$n\le300$ and all edge multiples $k\le15$, one finds:
for $r=2,3,4,7$ the systems contain $4,5,7,12$ polygons respectively, in lattices
of rank $2,3,4,3$, with $16,12,14,106$ vertex combinations; whereas for
$2\le r\le15$ with $r\notin\{2,3,4,7\}$ the only polygons in the span are the two
triangles defining the system, and no vertex closes. The four rich values are
exactly those for which $\cosh s$ has degree at most three over $\QQ$.
\end{observation}

\begin{remark}
The four fields are $\QQ(\sqrt[3]{2})$, $\QQ(\sqrt{33})$, $\QQ(\sqrt{33})$ and
$\QQ(\sqrt5)$. Only $r=3$ and $r=4$ share one, which is what forces their
difference to be rational; see Section~\ref{sec:limits}.
\end{remark}

\section{Limitations and open problems}\label{sec:limits}

\subsection*{Scaling the decision procedure}
The elimination of Theorem~\ref{thm:algebraic} was carried out with nested
radicals for the constants $\cos(\pi/n)$. This is adequate for small $n$ but
degenerates quickly. Of the four systems that resisted a first attempt, the two
with moderate coefficient fields have since been completed and appear in
Table~\ref{tab:edges}: $1.1555$, carrying pentagons and dodecagons, and $1.2537$,
which is the regular tiling $\{4,5\}$ with decagons and $30$-gons intruding. The
two that remain are $1.0377$, whose combination $\vc{3.4.8.40}$ carries a
$40$-gon, and $1.0531$, whose two combinations $\vc{3.5.6.18}$ and
$\vc{3.6^{2}.9}$ carry an $18$-gon together with a $9$-gon. Working throughout in
a single cyclotomic field $\QQ(\zeta_N)$ with $N=\operatorname{lcm}\{2n\}$, and
using modular or evaluation-interpolation resultants, should remove the
obstruction. The system at $1.0531$ is the one of greatest interest, since its
pentagons admit no symmetry and its smallest tilings are $16$-uniform.

\subsection*{Completeness of the catalogue}
The catalogue is believed complete, with the explicit caveat that nothing rules
out a system hiding between two listed edges. Theorem~\ref{thm:decide} decides
equality and inequality of two given edges; what is needed for completeness is a
bound, in terms of the degree $M$ of a vertex combination, on the height of
$P_m$, together with an enumeration of combinations up to that degree. Both
appear to be within reach.

\subsection*{Intrusion into a regular tiling}
Two entries of Table~\ref{tab:edges} have the same shape: a regular tiling
together with one further combination that brings in a pair of much larger
polygons. At $1.0612$ the tiling $\{5,4\}=\vc{5^{4}}$ is accompanied by
$\vc{3.4.10.20}$, and at $1.2537$ the tiling $\{4,5\}=\vc{4^{5}}$ is accompanied
by $\vc{4^{2}.10.30}$. Both intruding combinations contain a decagon, and in both
the second large polygon is a multiple of it. The two regular tilings are dual to
one another. Whether these are the first two members of a family, and what the
rule for the second large polygon is, we do not know.

\subsection*{Finiteness of a scaled hybrid system}
Observation~\ref{obs:106} presupposes the list of twelve polygons. We verified by
search that no other pair $(n,k)$ with $n\le2000$, $k\le60$ has its angle in the
rank-three lattice, but the argument is circular in the following sense: the
lattice was read off from the twelve polygons. What is missing is an argument
that no disjoint cluster of further polygons can close vertices among themselves
at the same edge.

\subsection*{Rank as an invariant}
In each system studied the polygons split into a part lying in
$\QQ\pi+\ZZ a$ and a small family carrying the remaining generators and one
relation among them: $\{\mathrm{A}10,\mathrm{A}20\}$ with
$\mathrm{A}10+\mathrm{A}20=\pi+\mathrm{A}3$ in the system of
Section~\ref{sec:lattice}; the apeirogons
$\{\mathrm{B}\infty,\mathrm{D}\infty,\mathrm{F}\infty\}$
with $3\mathrm{B}\infty+\mathrm{D}\infty+\mathrm{F}\infty=\pi$ in the system
$r=4$; and $\{\mathrm{A}6,\mathrm{B}3\}$ with $2\mathrm{A}6+\mathrm{B}3=\pi$ in
$r=3$. Whether this is a general phenomenon, and what governs the rank, we do not
know.

\subsection*{The shared field at \texorpdfstring{$r=3$}{r=3} and \texorpdfstring{$r=4$}{r=4}}
The identity $\cosh s_3-\cosh s_4=\tfrac12$ of Theorem~\ref{thm:rvalues} is a
consequence of both edges lying in $\QQ(\sqrt{33})$ with the same coefficient on
$\sqrt{33}$; their difference is then forced to be rational. Why these two
systems share a quadratic field, when no other pair in the family does, we cannot
say.

\section*{Reproducibility}

All computations are reproduced by scripts depending only on the
\texttt{mpmath} and \texttt{sympy} libraries. They recompute every closed form,
re-derive the enumerations both on the integer lattice and by direct
verification of the angle sums at $120$ decimal digits, and re-run the
completeness scans; a pass or failure is reported for each individual claim in
this paper.

\section*{Acknowledgements}

\emph{To be written. Authorship and order are to be settled by the authors; this
draft is circulated for comment. Computational work was carried out with the
assistance of Claude (Anthropic).}

\begin{thebibliography}{9}
\bibitem{abrate}
M. Abrate, S. Barbero, U. Cerruti and N. Murru,
\emph{Writing $\pi$ as sum of arctangents with linear recurrent sequences, Golden
mean and Lucas numbers}, arXiv:1409.6455 (2014).

\bibitem{calcutgauss}
J. S. Calcut,
\emph{Gaussian integers and arctangent identities for $\pi$},
Amer. Math. Monthly \textbf{116} (2009), 515--530.

\bibitem{calcutrational}
J. S. Calcut,
\emph{Rational angled hyperbolic polygons}, arXiv:1412.3825 (2014).

\bibitem{conwayjones}
J. H. Conway and A. J. Jones,
\emph{Trigonometric diophantine equations (On vanishing sums of roots of unity)},
Acta Arith. \textbf{30} (1976), 229--240.

\bibitem{ge}
G. Ge,
\emph{Algorithms related to multiplicative representations of algebraic numbers},
Ph.D. thesis, University of California, Berkeley, 1993.

\bibitem{maclachlanreid}
C. Maclachlan and A. W. Reid,
\emph{The Arithmetic of Hyperbolic 3-Manifolds},
Graduate Texts in Mathematics 219, Springer, 2003.

\bibitem{maiti}
A. Maiti,
\emph{On tilings of the hyperbolic plane by regular polygons}, arXiv:2302.05661.

\bibitem{pech}
P. Pech,
\emph{On the need of radical ideals in automatic proving: a theorem about
regular polygons}, and \emph{Computations of the area and radius of cyclic
polygons given by the lengths of sides}, in Automated Deduction in Geometry,
LNCS 3763, Springer (2006), 44--58.

\bibitem{robbins}
D. P. Robbins,
\emph{Areas of polygons inscribed in a circle},
Discrete Comput. Geom. \textbf{12} (1994), 223--236.

\bibitem{takeuchi}
K. Takeuchi,
\emph{Arithmetic triangle groups},
J. Math. Soc. Japan \textbf{29} (1977), 91--106.

\bibitem{zheng}
T. Zheng,
\emph{Computing multiplicative relations between roots of a polynomial},
arXiv:1912.07202 (2019).
\end{thebibliography}

\end{document}
